\documentclass[10pt]{article}
\usepackage[T1]{fontenc}
\usepackage[utf8]{inputenc}
\usepackage[french]{babel}
\usepackage{amsmath,amssymb,mathrsfs,mathpazo}
\usepackage{geometry}
\geometry{a4paper,margin=2cm}
\usepackage{graphicx,array,tabularx,tabularray,colortbl,multirow,multicol,booktabs,enumitem,pifont,ulem,fancybox,diagbox,cancel}
\usepackage{pstricks,pstricks-add,pst-plot,pst-tree,pst-node,pst-text,pst-func,pst-math,pst-bspline,pst-3dplot,pst-eucl,pst-solides3d}
\let\pscalculate\undefined
\usepackage{pst-bezier}
\usepackage{tikz,pgfplots,tkz-tab}
\usetikzlibrary{arrows,arrows.meta,patterns,calc,positioning,trees,decorations.markings,shapes.geometric,shapes.misc,shapes.arrows,intersections,decorations.pathmorphing,decorations.pathreplacing,backgrounds,shadows,babel}
\pgfplotsset{compat=1.18}
\usepgfplotslibrary{fillbetween}
\usepackage[np]{numprint}
\usepackage{listings,makecell,pas-tableur,twoopt}
\let\barre\undefined
\usepackage[pstricks]{tablvar}
\definecolor{darkgreen}{RGB}{0,100,0}
\providecommand{\red}{\color{red}}
\providecommand{\blue}{\color{blue}}
\providecommand{\green}{\color{green}}
\providecommand{\black}{\color{black}}
\providecommand{\white}{\color{white}}
\providecommand{\gray}{\color{gray}}
\providecommand{\cyan}{\color{cyan}}
\providecommand{\magenta}{\color{magenta}}
\providecommand{\yellow}{\color{yellow}}
\providecommand{\violet}{\color{violet}}
\providecommand{\orange}{\color{orange}}
\providecommand{\euro}{\texteuro}
\providecommand{\EUR}{\texteuro}
\providecommand{\ieme}{\textsuperscript{e}}
\providecommand{\degres}{\ensuremath{^\circ}}
\providecommand{\degree}{\ensuremath{^\circ}}
\providecommand{\vv}[1]{\overrightarrow{#1}}
\providecommand{\Coord}[1]{#1}
\providecommand{\decosix}{}
\providecommand{\starredbullet}{\ensuremath{\bullet}}
\pagestyle{empty}
\setlength{\parindent}{0pt}
\usepackage{adjustbox}\DeclareUnicodeCharacter{2236}{:}\DeclareUnicodeCharacter{2190}{\ensuremath{\leftarrow}}
\providecommand{\R}{\mathbb{R}}
\providecommand{\N}{\mathbb{N}}
\providecommand{\Z}{\mathbb{Z}}
\providecommand{\Q}{\mathbb{Q}}
\providecommand{\D}{\mathbb{D}}
\providecommand{\C}{\mathbb{C}}
\providecommand{\e}{{\mathrm{e}}}
\providecommand{\vect}[1]{\overrightarrow{#1}}
\providecommand{\vectt}[1]{\overrightarrow{\mathrm{#1}}}
\providecommand{\vvt}[1]{\overrightarrow{\mathrm{#1}}}
\providecommand{\np}[2][]{{#2~#1}}
\providecommand{\vv}[1]{{\overrightarrow{#1}}}
\renewcommand{\d}{{\mathrm{d}}}
\providecommand{\degres}{{^\circ}}
\providecommand{\degre}{{^\circ}}
\providecommand{\blue}{\color{blue}}
\providecommand{\red}{\color{red}}
\providecommand{\enskip}{\quad}
\providecommand{\llbracket}{\lbrack\!\lbrack}
\providecommand{\rrbracket}{\rbrack\!\rbrack}
\providecommand{\euro}{{\text{€}}}
\providecommand{\no}{{\text{n°}}}
\providecommand{\interval}[2]{{[#1;#2]}}
\providecommand{\intervFO}[2]{{[#1;#2[}}
\providecommand{\dotfill}{\ldots}

\providecommand{\ds}{\displaystyle}
\providecommand{\textding}[1]{#1}
\providecommand{\Oij}{\(\mathrm{O};\vec{\imath},\vec{\jmath}\)}
\providecommand{\Oijk}{\(\mathrm{O};\vec{\imath},\vec{\jmath},\vec{k}\)}
\providecommand{\Ouv}{\(\mathrm{O};\vec{u},\vec{v}\)}
\providecommand{\intervFF}[2]{\ensuremath{[#1;#2]}}
\renewcommand{\euro}{\ifmmode\text{\texteuro}\else\texteuro\fi}
\DeclareUnicodeCharacter{00B0}{\ifmmode^{\circ}\else\textdegree\fi}
\DeclareUnicodeCharacter{20AC}{\ifmmode\text{\texteuro}\else\texteuro\fi}
% Énoncés : Éducation nationale. Transcription et corrigés : APMEP, Denis Vergès et contributeurs. Découpage et classement : CoopMaths, Sébastien Lozano, Éric Elter et contributeurs. Adaptation de la présentation : Maths Explorer.
% CC BY-SA : https://coopmaths.fr/www/about/
% https://coopmaths.fr/alea/static/bac/2026/tex/bac_2026_05_sujet2_ameriquenord_4_cor.tex
\begin{document}
\noindent\textit{APMEP / CoopMaths. CC BY-SA. Présentation adaptée par Maths Explorer.}\par\medskip
On considère la fonction $f$ définie sur l'intervalle $]0~;~+\infty[$ par
$f(x) = x(\ln x)^2$.





\begin{enumerate}
	\item On cherche la limite de la fonction $f$ en $+\infty$.

	Pour tout réel $x$ strictement positif, $f(x)=x\left(\ln x\right)^2=x \times \ln(x) \times \ln(x)$. 
	
	$\begin{cases}
		\lim\limits_{x \to +\infty} \ln(x)= +\infty\\
		\lim\limits_{x\to +\infty}  x = +\infty
	\end{cases}$ donc, par produit,  $\lim\limits_{x \to +\infty} f(x)= +\infty$.
	
	\item Pour tout réel x > 0, on pose $g(x) = x\ln x$.
	\begin{enumerate}
		\item Soit $x \in ]0~;~+\infty[$, 

		$4\left(g\left(\sqrt x\right)\right) ^2= 4 \left( \sqrt{x} \ln\big(\sqrt{x}\big) \right) ^2= 4 \left( \sqrt{x} \times \dfrac{1}{2} \ln(x)\right)^2 = 4 \times x \times \dfrac{1}{4} \left( \ln x\right) ^2 = x \left( \ln x\right) ^2 =f(x)$

On a bien :  pour tout réel $x > 0$, on a $f(x) = 4\left(g\left(\sqrt x\right)\right)$.
		\item On cherche $\displaystyle\lim_{x \to 0} f(x)$.

$\begin{cases}
\lim\limits_{x \to 0} \sqrt{x} =0 \\
\lim\limits_{X \to 0} X\ln X = 0 \, \text{par croissances comparées} \\
\end{cases}$ donc, par composition, $\lim\limits_{x\to 0} g(\sqrt x)=0$

D'où, par produit, $\lim\limits_{x\to 0} f(x)=0$. 
	\end{enumerate}
\item On étudie les variations de la fonction $f$ sur l'intervalle $]0~;~+\infty[$.
	\begin{enumerate}
		\item $f$ est de la forme $uv$ avec : \quad $\forall x \in ]0;+\infty[  \begin{cases}
u(x)=x\\ v(x)=(\ln x)^2\end{cases}$

On a alors :\quad $\forall x \in ]0;+\infty[ 
\begin{cases} u'(x)=1 \\ v'(x)=2 \times \dfrac{1}{x} \times \ln x 
\end{cases}$ 

Or :\quad $f' = u'v+v'u$ 

Donc, $\begin{aligned}[t]\forall x \in ] 0 ;+\infty[, \quad f'(x)&= 1 \times (\ln x)^2 + x \times \dfrac{2 \ln x}{x}\\
		&= \ln x \times \ln x + 2 \ln x \\
		&=(\ln x)(\ln x +2)
\end{aligned}$
		
Donc sur l'intervalle $]0~;~+\infty[$, \: $f'(x) = (\ln x)(2 + \ln x)$.

	\item On en déduit les variations de la fonction $f$ sur l'intervalle $]0~;~+\infty[$.

Étudions le signe de $f'(x)$ : 

$2+ \ln(x)> 0 \Longleftrightarrow \ln x > -2 \Longleftrightarrow x > \e^{-2}$ 

$\ln x > 0 \Longleftrightarrow x >1$ 

On a donc le tableau de variations suivant : 

\hfill~\begin{tikzpicture}
	\tkzTabInit[lgt=3,espcl=2]{$x$/0.7,
	signe de $\ln x$/0.7,
	signe de $2 + \ln x$/0.7,
	signe de $f'(x)$/0.7,
	variations de $f$/1.5}{$0$,$\e^{-2}$,1,$+\infty$}
	\tkzTabLine{,-,t,-,z,+,}
	\tkzTabLine{,-,z,+,t,+,}
	\tkzTabLine{,+,z,-,z,+,}
	\tkzTabVar{-/0,+/$f\left(\e^{-2}\right)$,-/0,+/{$+\infty$}}
		\end{tikzpicture}\hfill~

		\item $f\left(\e^{-2}\right) = \e^{-2} \left( \ln \left(\e^{-2}\right) \right) ^2= \e^{-2} \times (-2)^2= 4 \e^{-2}$

Le maximum de la fonction $f$ sur l'intervalle $]0~;~1]$ est $4\e^{-2}$.
	\end{enumerate}

\item 
	\begin{enumerate}
		\item Sur l'intervalle $]0~;~1]$, le maximum est $4\e^{-2} \approx 0,54 < 2$ donc l'équation $f(x) = 2$ n'admet pas de solution sur cet intervalle.

Sur l'intervalle $[1; +\infty[$ : 
\begin{itemize}
	\item $f$ est continue car dérivable ;
	\item $f$ est strictement croissante ;
	\item $f(1) =0 <2$ et $\lim\limits_{x \to +\infty} f(x)= +\infty >2$ 
\end{itemize}

Donc, d'après le théorème de la bijection, l'équation $f(x)=2$ admet une unique solution sur l'intervalle $[1~;~+\infty[$.

donc l'équation admet une unique solution  $\alpha$ sur l'intervalle $]0~;~+\infty[$.

\item À l'aide de la calculatrice, on trouve $\alpha \approx 2,46$ donc : \quad  $ 2,4 < \alpha <2,5$. 
	\end{enumerate}
\item Soit $a$ un nombre réel appartenant à l'intervalle ]0~;~1].
	\begin{enumerate}
		\item La fonction $f$ est positive sur l'intervalle $[a;1]$ donc $\displaystyle\int_a^1 f(x)\: \text{d}x$ est égal à l'aire, en unité d'aire,  du domaine délimité par les droites d'équation $x=a$ et $x=1$, par l'axe des abscisses et la courbe. 
		
		\item  $x \longmapsto x \left(\ln x \right)^2$\quad est de la forme $u' v$ avec \quad $\begin{cases}u'(x)=x\\ v(x)=(\ln x)^2 \end{cases}$

Ce qui implique : \quad $\begin{cases}u(x)=\dfrac{1}{2}x^2\\v'(x)=2 \times \dfrac{1}{x}\ln (x)\end{cases}$

Les fonctions $u$ et $v$ sont continues et dérivables sur $[a~;~1]$ et les fonctions $u'$ et $v'$ sont continues sur $[a~;~1]$, donc , d'après la formule d'intégration par parties :

$\displaystyle \int_{a}^{1} u'(x)v(x)\, \mathrm{d}x = \Big [ u(x)v(x)\strut\Big]_{a}^{1} - \int_{a}^{1} u(x)v'(x) \, \mathrm{d}x$\

donc 
$\begin{aligned}[t]
\displaystyle\int_a^1 f(x)\, \mathrm{d}x  & =\Big[ \dfrac{x^2}{2 }\times (\ln x)^2 \Big]_a^{1}-\int_a^{1} \dfrac{x^2}{2} \times \dfrac{2\ln x}{x} \, \mathrm{d}x \\
	& = \dfrac{1^2}{2} (\ln 1)^2  - \dfrac{a^2}{2} (\ln a)^2  - \int_a^{1} x \ln(x) \, \mathrm{d}x \\
	&= - \dfrac{a^2}{2} (\ln a)^2  - \int_a^{1} x \ln(x) \, \mathrm{d}x
\end{aligned}	$

Ce qui est le résultat demandé. 

	\item  $x \longmapsto x \ln x $\quad est de la forme $u' v$ avec \quad $\begin{cases}u'(x)=x\\ v(x)=\ln x \end{cases}$

Ce qui implique : \quad $\begin{cases}u(x)=\dfrac{1}{2}x^2\\v'(x)=\dfrac{1}{x}\end{cases}$

Les fonctions $u$ et $v$ sont continues et dérivables sur $[a~;~1]$ et les fonctions $u'$ et $v'$ sont continues sur $[a~;~1]$, donc , d'après la formule d'intégration par parties :

$\displaystyle \int_{a}^{1} u'(x)v(x)\, \mathrm{d}x = \Big [ u(x)v(x)\strut\Big]_{a}^{1} - \int_{a}^{1} u(x)v'(x) \, \mathrm{d}x$\

donc 
$\begin{aligned}[t]
\displaystyle\int_a^1 x \ln (x)\, \mathrm{d}x  & =\Big[ \dfrac{x^2}{2 }\times \ln (x)\Big]_a^{1}-\int_a^{1} \dfrac{x^2}{2} \times \dfrac{1}{x} \, \mathrm{d}x \\
	& =  \dfrac{1^2}{2} \ln (1)  - \dfrac{a^2}{2} \ln (a)  - \int_a^{1} \dfrac{1}{2}x \, \mathrm{d}x \\
	&= - \dfrac{a^2}{2} \ln (a)  - \Big[ \dfrac{1}{2} \times \dfrac{x^2}{2 }\Big]_a^{1}\\
	&= - \dfrac{a^2}{2} \ln (a)  - \Big[ \dfrac{1}{2} \times \dfrac{1^2}{2 } - \dfrac{1}{2} \times \dfrac{a^2}{2 }  \Big]\\
	&= - \dfrac{a^2}{2} \ln (a)  -  \dfrac{1}{4} +  \dfrac{a^2}{4 } 
\end{aligned}$

D'où :
$\begin{aligned}[t]
\displaystyle\int_a^1 f(x)\: \text{d}x & =  - \dfrac{a^2}{2} (\ln a)^2  - \int_a^{1} x \ln(x) \, \mathrm{d}x \\
	&= - \dfrac{a^2}{2}(\ln a)^2  - \left( - \dfrac{a^2}{2} \ln (a)  -  \dfrac{1}{4} +  \dfrac{a^2}{4 }  \right) \\
	&= - \dfrac{a^2}{2}(\ln a)^2 + \dfrac{a^2}{2} \ln a + \dfrac14 - \dfrac{a^2}{4}
		\end{aligned}$

Ce qui est le résultat demandé.
\item  On cherche  la limite de $\displaystyle\int_a^1 f(x)\: \text{d}x$ quand $a$ tend vers 0.

$\displaystyle\int_a^1 f(x)\: \text{d}x = - \dfrac{a^2}{2}(\ln a)^2 + \dfrac{a^2}{2} \ln a + \dfrac14 - \dfrac{a^2}{4} = - \dfrac{1}{2} \left(a \ln a\right)^2+ \dfrac{a}{2} \times  a \ln a + \dfrac{1}{4} - \dfrac{a^2}{4}$.

Par croissances comparées : $\lim\limits_{a \to 0} a \ln a = 0$ 

donc : par produit $\lim\limits_{a \to 0} \left(a \ln a \right)^2= 0$ et $\lim\limits_{a \to 0} \dfrac{a}{2} \times a\ln a=0$.

De plus, $\lim\limits_{a \to 0}  \dfrac{1}{4} - \dfrac{a^2}{4}=\dfrac{1}{4}$ 

donc, par somme, $\lim\limits_{a \to 0 }	\displaystyle\int_a^1 f(x)\: \text{d}x  = \dfrac{1}{4}$

	\end{enumerate}
\end{enumerate}

\end{document}
